\chapter{上下文压缩与 Token 管理}

\section{问题背景}

LLM 的上下文窗口是有限的。Claude 的上下文窗口约为 200K tokens，但在长时间的编程会话中，消息历史可以轻松超过这个限制。Claude Code 设计了一套多层压缩系统来解决这个问题。

\section{五层压缩管线}

\begin{lstlisting}[language=TypeScript,breaklines=true]
原始消息历史
    │
    ▼
① 工具结果预算 (Tool Result Budget)
    │  控制单个工具输出的大小
    ▼
② Snip 压缩 (History Snip)
    │  移除旧的对话片段
    ▼
③ 微压缩 (Microcompact)
    │  压缩单个工具输出的冗余
    ▼
④ 上下文折叠 (Context Collapse)
    │  多轮交互折叠为摘要
    ▼
⑤ 自动压缩 (Auto Compact)
    │  接近窗口限制时的全局压缩
    ▼
发送到 API
\end{lstlisting}

\section{第一层：工具结果预算}

\begin{lstlisting}[language=TypeScript,breaklines=true]
messagesForQuery = await applyToolResultBudget(
  messagesForQuery,
  toolUseContext.contentReplacementState,
  persistReplacements ? records =>
    void recordContentReplacement(records, toolUseContext.agentId).catch(logError)
    : undefined,
  new Set(
    toolUseContext.options.tools
      .filter(t => !Number.isFinite(t.maxResultSizeChars))
      .map(t => t.name),
  ),
)
\end{lstlisting}

\textbf{工作原理}：
\begin{itemize}[nosep]
  \item 每个工具可以定义 \code{maxResultSizeChars} 限制
  \item 超出限制的工具结果被替换为摘要
  \item 替换决策通过 \code{contentReplacementState} 持久化
  \item 支持会话恢复时重建替换状态
\end{itemize}

\textbf{与微压缩的关系}：工具结果预算在微压缩之前运行。由于缓存微压缩完全通过 \code{tool\_use\_id} 操作（不检查内容），内容替换对它是透明的，两者可以干净地组合。

\section{第二层：Snip 压缩}

\begin{lstlisting}[language=TypeScript,breaklines=true]
if (feature('HISTORY_SNIP')) {
  const snipResult = snipModule!.snipCompactIfNeeded(messagesForQuery)
  messagesForQuery = snipResult.messages
  snipTokensFreed = snipResult.tokensFreed
  if (snipResult.boundaryMessage) {
    yield snipResult.boundaryMessage
  }
}
\end{lstlisting}

Snip 压缩是一种基于"片段"的压缩方式：
\begin{itemize}[nosep]
  \item 标识旧的对话片段
  \item 移除不再需要的旧片段
  \item 通过 \code{snipReplay} 回调处理 snip 边界
  \item 释放的 Token 数量传递给自动压缩阈值计算
\end{itemize}

\subsection{Snip 边界处理}

\begin{lstlisting}[language=TypeScript,breaklines=true]
// QueryEngine 配置中的 snipReplay 回调
snipReplay?: (
  yieldedSystemMsg: Message,
  store: Message[],
) => { messages: Message[]; executed: boolean } | undefined
\end{lstlisting}

snipReplay 被注入到 QueryEngine 中，使得 Feature Flag 控制的字符串保持在门控模块内，QueryEngine 保持可测试。

\section{第三层：微压缩}

\begin{lstlisting}[language=TypeScript,breaklines=true]
const microcompactResult = await deps.microcompact(
  messagesForQuery,
  toolUseContext,
  querySource,
)
messagesForQuery = microcompactResult.messages
\end{lstlisting}

微压缩处理单个工具输出中的冗余信息，例如：
\begin{itemize}[nosep]
  \item 重复的文件列表
  \item 冗长的命令输出
  \item 重复的错误信息
\end{itemize}

\subsection{缓存微压缩}

\begin{lstlisting}[language=TypeScript,breaklines=true]
const pendingCacheEdits = feature('CACHED_MICROCOMPACT')
  ? microcompactResult.compactionInfo?.pendingCacheEdits
  : undefined
\end{lstlisting}

缓存微压缩是一种优化：使用 Anthropic API 的缓存编辑功能，延迟边界消息直到 API 响应后才发送（使用实际的 \code{cache\_deleted\_input\_tokens} 而非估计值）。

\section{第四层：上下文折叠}

\begin{lstlisting}[language=TypeScript,breaklines=true]
if (feature('CONTEXT_COLLAPSE') && contextCollapse) {
  const collapseResult = await contextCollapse.applyCollapsesIfNeeded(
    messagesForQuery,
    toolUseContext,
    querySource,
  )
  messagesForQuery = collapseResult.messages
}
\end{lstlisting}

\textbf{设计特点}：
\begin{itemize}[nosep]
  \item 上下文折叠在自动压缩之前运行
  \item 如果折叠使 Token 数降到自动压缩阈值以下，自动压缩就不会触发
  \item 折叠视图是\textbf{读时投影}，不修改 REPL 的完整历史
  \item 摘要消息存储在折叠存储中，而非 REPL 消息数组
  \item \code{projectView()} 在每次查询入口重放提交日志
\end{itemize}

\section{第五层：自动压缩}

\begin{lstlisting}[language=TypeScript,breaklines=true]
import {
  calculateTokenWarningState,
  isAutoCompactEnabled,
  type AutoCompactTrackingState,
} from './services/compact/autoCompact.js'

import { buildPostCompactMessages } from './services/compact/compact.js'
\end{lstlisting}

自动压缩是最后的安全网：
\begin{itemize}[nosep]
  \item 当 Token 使用接近上下文窗口限制时自动触发
  \item 使用 LLM 生成整个对话的摘要
  \item 替换原有消息历史为摘要 + 最近的消息
  \item 通过 \code{/compact} 命令也可以手动触发
\end{itemize}

\subsection{响应式压缩}

\begin{lstlisting}[language=TypeScript,breaklines=true]
const reactiveCompact = feature('REACTIVE_COMPACT')
  ? require('./services/compact/reactiveCompact.js')
  : null
\end{lstlisting}

响应式压缩是自动压缩的一个变种——当 API 返回 "prompt too long" 错误时触发，是对主动压缩的被动补充。

\section{Token 计数与估算}

\begin{lstlisting}[language=TypeScript,breaklines=true]
import {
  doesMostRecentAssistantMessageExceed200k,
  finalContextTokensFromLastResponse,
  tokenCountWithEstimation,
} from './utils/tokens.js'
\end{lstlisting}

Token 计数使用两种策略：
\begin{enumerate}[nosep]
  \item \textbf{精确计数}：从 API 响应的 \code{usage} 字段获取
  \item \textbf{估算}：在 API 调用之前，使用估算方法预测 Token 数量
\end{enumerate}

\section{Token 预算系统}

\begin{lstlisting}[language=TypeScript,breaklines=true]
const budgetTracker = feature('TOKEN_BUDGET')
  ? createBudgetTracker()
  : null
\end{lstlisting}

Token 预算系统（\code{TOKEN\_BUDGET} Feature Flag）追踪每个回合的 Token 消耗：

\begin{lstlisting}[language=TypeScript,breaklines=true]
import {
  getCurrentTurnTokenBudget,
  getTurnOutputTokens,
  incrementBudgetContinuationCount,
} from './bootstrap/state.js'

import { createBudgetTracker, checkTokenBudget } from './query/tokenBudget.js'
\end{lstlisting}

\subsection{任务预算}

\begin{lstlisting}[language=TypeScript,breaklines=true]
// API 的 task_budget (output_config.task_budget)
// 与 tokenBudget 的 +500k 自动继续功能不同
// `total` 是整个代理回合的预算
// `remaining` 从累积 API 使用量中计算
taskBudget?: { total: number }
\end{lstlisting}

任务预算通过 API 的 \code{output\_config.task\_budget} 参数控制。压缩后需要重新计算剩余预算：

\begin{lstlisting}[language=TypeScript,breaklines=true]
// 跨压缩边界的 task_budget.remaining 追踪
// 未压缩时服务器可以看到完整历史并自行处理倒计时
// 压缩后服务器只看到摘要，会低估消耗
// remaining 告诉它压缩前的最终窗口
let taskBudgetRemaining: number | undefined = undefined
\end{lstlisting}

\section{max\_output\_tokens 恢复}

\begin{lstlisting}[language=TypeScript,breaklines=true]
const MAX_OUTPUT_TOKENS_RECOVERY_LIMIT = 3

function isWithheldMaxOutputTokens(msg) {
  return msg?.type === 'assistant' && msg.apiError === 'max_output_tokens'
}
\end{lstlisting}

当模型输出被 \code{max\_output\_tokens} 截断时：
\begin{enumerate}[nosep]
  \item 错误消息被"扣留"，不发送给 SDK 调用方
  \item 自动向 API 发送继续请求
  \item 最多重试 3 次
  \item 只有确认恢复失败后才将错误报告给调用方
\end{enumerate}

\section{成本追踪}

\begin{lstlisting}[language=TypeScript,breaklines=true]
// cost-tracker.ts
export function addToTotalSessionCost(
  cost: number,
  usage: Usage,
  model: string,
): number {
  const modelUsage = addToTotalModelUsage(cost, usage, model)
  addToTotalCostState(cost, modelUsage, model)

  // OpenTelemetry 计数器
  getCostCounter()?.add(cost, attrs)
  getTokenCounter()?.add(usage.input_tokens, { ...attrs, type: 'input' })
  getTokenCounter()?.add(usage.output_tokens, { ...attrs, type: 'output' })
  getTokenCounter()?.add(usage.cache_read_input_tokens ?? 0, {
    ...attrs, type: 'cacheRead'
  })
  getTokenCounter()?.add(usage.cache_creation_input_tokens ?? 0, {
    ...attrs, type: 'cacheCreation'
  })

  // 处理 advisor 使用量
  for (const advisorUsage of getAdvisorUsage(usage)) {
    const advisorCost = calculateUSDCost(advisorUsage.model, advisorUsage)
    totalCost += addToTotalSessionCost(advisorCost, advisorUsage, advisorUsage.model)
  }
  return totalCost
}
\end{lstlisting}

成本追踪记录：
\begin{itemize}[nosep]
  \item 每个模型的输入/输出 Token
  \item 缓存读取/创建 Token
  \item 网页搜索请求数
  \item Advisor（辅助模型）的使用量
  \item USD 成本计算
\end{itemize}

\subsection{会话成本持久化}

\begin{lstlisting}[language=TypeScript,breaklines=true]
export function saveCurrentSessionCosts(fpsMetrics?: FpsMetrics): void {
  saveCurrentProjectConfig(current => ({
    ...current,
    lastCost: getTotalCostUSD(),
    lastAPIDuration: getTotalAPIDuration(),
    lastTotalInputTokens: getTotalInputTokens(),
    lastTotalOutputTokens: getTotalOutputTokens(),
    lastSessionId: getSessionId(),
    // ...
  }))
}
\end{lstlisting}

成本在会话切换前保存到项目配置中，支持跨会话的成本恢复。

\section{设计启示}

\subsection{多层防御}

五层压缩不是冗余设计，每层解决不同粒度的问题。这种"洋葱式"架构确保了即使某一层未能充分压缩，后续层仍能处理。

\subsection{缓存感知}

压缩系统与 Prompt Cache 深度集成。例如，缓存微压缩利用缓存编辑 API 避免不必要的缓存失效。

\subsection{预算管理}

Token 预算和成本追踪的细粒度设计，使得企业用户可以精确控制 AI 的资源消耗。

\bigskip\noindent\rule{\textwidth}{0.4pt}\bigskip

\textbf{下一章}：\href{./chapter-10-coordinator.md}{多智能体协调 (Coordinator Mode)}
